\subsection{有理整数的整除性}

正如第 1 节所指出的，有理整数环 Z 是一个主理想整环；其分式域为 Q。Z 的可逆元乘法群 U 包含两个元素 1 和 -1。群 \(Q_{+}^{*}\) （由 \textgreater0 的有理数组成）恰好包含来自 Q 的每个相伴元素类中的一个元素；因此它与 Q 的主分式理想组成的乘法群 \(P^{*} = Q^{*}/U\) 同构，通常把它们等同看待。特别地，当在域 Q 中（相对于环 Z）使用 gcd 或 lcm 时，应理解为这些是 \(\geqslant 0\) 的元素；这一约定使我们能够谈论一族有理数的最大公因子和最小公倍元。

Z 中的不可约整数 \textgreater0 正是我们称之为素数的那些数（I，第 50 页）（它们有时被称为有理素数）；Z 的每个不可约元素因此具有 p 或 -p 的形式，其中 p 是素数，且素数集合 P 是 Z 的不可约元素的一个代表系。

命题5. --- 素数集合是无限的。

事实上，给定任意有限族 \((p_{i})\) \((1 \leqslant i \leqslant n)\) （由互不相同的素数组成），数 \(\left(\prod_{i=1}^{n} p_{i}\right) + 1\) （它是 \textgreater1 的）的任何素因子 q 与所有 \(p_{i}\) 不同，否则它将整除 1。

